\documentclass[11pt]{article}
\usepackage[a4paper,margin=28mm]{geometry}
\usepackage{amsmath,amssymb,amsthm}
\usepackage[colorlinks=true,linkcolor=blue,urlcolor=blue,citecolor=blue]{hyperref}
\newtheorem{theorem}{Theorem}
\newtheorem{lemma}{Lemma}
\newcommand{\Q}{\mathbb Q}
\newcommand{\Z}{\mathbb Z}
\newcommand{\F}{\mathbb F}
\DeclareMathOperator{\End}{End}
\DeclareMathOperator{\Gal}{Gal}
\DeclareMathOperator{\GL}{GL}
\DeclareMathOperator{\PGL}{PGL}
\title{An obstruction to rational \(35\)-torsion on geometrically split
abelian surfaces}
\author{}
\date{20 October 2026}

\begin{document}
\maketitle

\begin{theorem}\label{thm:main}
Let \(A/\Q\) be an abelian surface that is simple over \(\Q\) and
geometrically split. Then \(A(\Q)\) has no point of order \(35\).
\end{theorem}

Here ``geometrically split'' means that
\(A_{\overline{\Q}}\) is isogenous to a product of elliptic curves.
In particular, the groups \([35]\) and \([70]\) cannot occur as the
rational torsion subgroup of a geometrically split genus-\(2\)
Jacobian that is simple over \(\Q\).

The hypothesis of simplicity over \(\Q\) is essential. A product
of elliptic curves over \(\Q\) with rational points of orders \(5\)
and \(7\) has a rational point of order \(35\).

\section{Inputs on elliptic torsion and reduction}

By Derickx--Najman \cite[Theorem 1.1]{DN}, no elliptic curve over
a quartic number field has a point of order \(35\). The same holds
over fields of degree \(1\) or \(2\), by extending the ground field
to degree \(4\).

We also use Clark's theorem \cite[Theorem 5]{Clark}: if an abelian
surface over \(\Q\) has everywhere potentially good reduction, the
order of its rational torsion subgroup belongs to
\[
 \{1,\ldots,16,18,19,20,22,24,25,28,30,36,48,60,72\}.
\]
None of these integers is divisible by \(35\).

\section{Combining torsion primes in an elliptic isogeny class}

\begin{lemma}\label{lem:combine}
Let \(p\) and \(q\) be distinct primes. Suppose \(E_p\) and \(E_q\)
are \(K\)-isogenous elliptic curves over a number field \(K\), with
points of orders \(p\) and \(q\) in \(E_p(K)\) and \(E_q(K)\),
respectively. Their \(K\)-isogeny class contains an elliptic curve
with a \(K\)-rational point of order \(pq\).
\end{lemma}

\begin{proof}
Choose a \(K\)-isogeny \(f:E_p\to E_q\), and let \(H\) be the
\(q\)-primary subgroup of \(\ker f\). Factor \(f\) as
\[
 E_p\xrightarrow{a}E'=E_p/H\xrightarrow{b}E_q.
\]
The degree of \(a\) is a power of \(q\), so \(a\) preserves the
order of the given \(p\)-torsion point. The degree of \(b\) is
prime to \(q\), so its dual maps the given \(q\)-torsion point
injectively into \(E'(K)\). The sum of the resulting points has
order \(pq\).
\end{proof}

In the rest of the proof, suppose that \(A(\Q)\) contains a point
of order \(35\). Choose nonzero points
\(P_p\in A(\Q)[p]\) for \(p\in\{5,7\}\).

\section{Nonisogenous factors and the CM case}

Suppose first that \(A_{\overline{\Q}}\) is isogenous to
\(E_1\times E_2\), where \(E_1\) and \(E_2\) are not
geometrically isogenous. There are exactly two geometric elliptic
subvarieties of \(A\), and Galois permutes them. Since \(A\) is
simple over \(\Q\), this action is nontrivial. Thus the two
subvarieties \(E\) and \(E^\sigma\) are defined over a quadratic
field \(K\), and its nontrivial automorphism \(\sigma\) interchanges
them.

Put \(F=A_K/E\). The restriction \(E^\sigma\to F\) of the quotient
map is a \(K\)-isogeny. Indeed, its kernel is finite, and its image
is a nonzero elliptic subvariety of \(F\). Conjugating, we find
that \(F^\sigma\) is \(K\)-isogenous to \(E\).

For each \(p\in\{5,7\}\), the image of \(P_p\) in \(F(K)\) either
has order \(p\) or is zero. In the first case \(F^\sigma(K)\)
contains a point of order \(p\). In the second case
\(P_p\in E(K)\). Thus, for each \(p\in\{5,7\}\), the \(K\)-isogeny
class of \(E\) contains a curve with a \(K\)-rational point of
order \(p\).

By Lemma~\ref{lem:combine}, this class contains an elliptic curve
over \(K\) with a point of order \(35\). This contradicts the
quartic torsion classification, after a quadratic extension of
\(K\). This argument also covers nonisogenous CM factors.

Suppose next that \(A_{\overline{\Q}}\) is isogenous to \(E^2\),
where \(E\) has complex multiplication. CM elliptic curves have
everywhere potentially good reduction. This property is preserved
by products and isogenies, and descends from a finite extension
of the ground field. Hence \(A\) has everywhere potentially good
reduction, and Clark's theorem gives a contradiction.

It remains to consider the case
\(A_{\overline{\Q}}\sim E^2\), with \(E\) non-CM.

\section{The endomorphism field}

In this case
\[
 D=\End^0(A_{\overline{\Q}})\simeq M_2(\Q).
\]
Let \(L/\Q\) be the smallest field over which all geometric
endomorphisms of \(A\) are defined, and put
\(\Gamma=\Gal(L/\Q)\). The extension \(L/\Q\) is finite Galois.
For a number field \(K\subset\overline{\Q}\), write
\(G_K=\Gal(\overline{\Q}/K)\).
By Skolem--Noether, the faithful action of \(\Gamma\) on \(D\)
identifies it with a finite subgroup of
\[
 \operatorname{Aut}_{\Q}(M_2(\Q))=\PGL_2(\Q).
\]

A finite subgroup of \(\PGL_2(\mathbb R)\) is cyclic or dihedral.
Indeed, its action on the hyperbolic plane fixes a point, and the
action on the tangent plane embeds it in \(O(2)\). If an element
of \(\PGL_2(\Q)\) has order \(n>2\), the ratio of the eigenvalues
of a representative \(M\) is a primitive \(n\)-th root of unity
\(\zeta\), and
\[
 \frac{(\operatorname{tr}M)^2}{\det M}
 =2+\zeta+\zeta^{-1}\in\Q.
\]
Hence \(n\in\{3,4,6\}\). It follows that
\[
 \Gamma\in\{C_1,C_2,C_3,C_4,C_6,D_4,D_6,D_8,D_{12}\},
\]
where the subscript on a dihedral group denotes its order.
In particular, \(D_4=C_2\times C_2\).

\section{An irreducible projective factor and the cyclotomic field}

\begin{lemma}\label{lem:cyclotomic}
Let \(p\) be an odd prime. Suppose \(\Gamma\) preserves a maximal
\(\Z_p\)-order \(\mathcal R\) in \(M_2(\Q_p)\), and the induced
projective action on its natural two-dimensional residual module
is irreducible over \(\F_p\). If \(A(\Q)[p]\neq0\), then
\(\Q(\zeta_p)\subseteq L\).
\end{lemma}

\begin{proof}
Write \(\mathcal R=\End_{\Z_p}(\Lambda)\) for a lattice
\(\Lambda\subset\Q_p^2\). Let \(T=T_p(A)\), and put
\[
 T'=\mathcal R T.
\]
Since \(\mathcal R\) is finite free over \(\Z_p\), this is a full
lattice in \(V_p(A)\). It is \(G_{\Q}\)-stable, since \(T\) and
\(\mathcal R\) are stable. After choosing a basis of \(\Lambda\),
the matrix units of \(\mathcal R=M_2(\Z_p)\) give
\[
 T'\simeq\Lambda\otimes_{\Z_p}W_0,
\]
where \(W_0=e_{11}T'\) is free of rank \(2\). The off-diagonal
matrix units identify \(e_{11}T'\) and \(e_{22}T'\), giving the
displayed decomposition.

Put \(\overline V=T'/pT'\). By independence of the semisimplified
residual representation from the choice of a stable lattice,
\(\overline V\) and \(T/pT=A[p]\) have the same composition
factors. Since \(A(\Q)[p]\neq0\), the module \(\overline V\) has a
trivial composition factor.

The action of \(G_{\Q}\) on \(\overline V\) normalizes
\(M_2(\F_p)\) acting on \(U=\Lambda/p\Lambda\). Its projective
action on \(U\) factors through \(\Gamma\). Pull back the central
extension
\[
 1\longrightarrow\F_p^\times\longrightarrow\GL_2(\F_p)
 \longrightarrow\PGL_2(\F_p)\longrightarrow1
\]
along this projective action of \(G_{\Q}\), and denote the
resulting group by \(\widetilde G\).
There are two-dimensional \(\widetilde G\)-modules \(U\) and \(W\)
with
\[
 \overline V\simeq U\otimes_{\F_p}W.
\]
Explicitly, an element \((g,B)\in\widetilde G\) acts on \(U\)
through \(B\). Multiplying the action of \(g\) on \(\overline V\)
by \(B^{-1}\) on the first factor gives an operator commuting
with \(M_2(\F_p)\); it therefore acts on the second factor and
defines the action on \(W\). All these actions factor through a
common finite quotient of \(\widetilde G\).

The module \(U\) is irreducible by assumption. The normal subgroup
\[
 H=\{(h,I):h\in G_L\}\subseteq\widetilde G
\]
acts trivially on \(U\). The trivial composition factor of
\(\overline V\) remains trivial on pullback to \(\widetilde G\).

If \(W\) were reducible, its two composition factors would be
characters. Tensoring a composition series with \(U\), the
composition factors of \(U\otimes W\) would then be irreducible
two-dimensional twists of \(U\). This contradicts the existence
of a trivial composition factor. Thus \(W\) is irreducible.

The restriction \(W|_H\) is semisimple. To see this in modular
characteristic, choose a simple \(H\)-submodule of \(W\). Its
\(\widetilde G\)-translates are simple \(H\)-submodules, by
normality of \(H\). Their sum is a nonzero
\(\widetilde G\)-submodule of \(W\), hence all of \(W\).

Restriction to \(H\) preserves the existence of a trivial
composition factor. Since \(\overline V|_H\) is two copies of
the semisimple module \(W|_H\), it follows that \(W^H\neq0\).
Normality of \(H\) makes \(W^H\) a \(\widetilde G\)-submodule.
Irreducibility gives \(W^H=W\). Hence \(\overline V|_{G_L}\)
is the trivial four-dimensional representation.

Choose a polarization of \(A\) over \(\Q\), which exists for every
abelian variety over \(\Q\). It gives an isomorphism of rational
\(p\)-adic representations
\[
 V_p(A)\simeq V_p(A)^*\otimes_{\Q_p}\Q_p(1).
\]
Transport the lattice
\(\operatorname{Hom}_{\Z_p}(T',\Z_p)\otimes_{\Z_p}\Z_p(1)\)
into \(V_p(A)\) by this isomorphism. Its reduction and
\(\overline V\) have the same semisimplification as
\(G_L\)-modules. The latter consists of four copies of the trivial
character, while the former consists of four copies of the
mod-\(p\) cyclotomic character \(\chi_p|_{G_L}\). Therefore
\(\chi_p|_{G_L}=1\), which is equivalent to
\(\Q(\zeta_p)\subseteq L\).
\end{proof}

The polarization is used on the rational Tate module. Thus the
argument applies even if its degree is divisible by \(p\). The
comparison of lattices uses only their residual
semisimplifications.

\section{The primes \(5\) and \(7\)}

Suppose first that \(\Gamma\) has a normal rotation subgroup
\(\langle r\rangle\) of order \(n=3\) or \(6\). This occurs for
\[
 \Gamma\in\{C_3,C_6,D_6,D_{12}\}.
\]
Normalize a matrix representative of \(r\) to have trace \(1\).
The resulting matrix \(R\) has characteristic polynomial
\[
 \begin{cases}
  X^2-X+1,&n=3,\\
  X^2-X+1/3,&n=6.
 \end{cases}
\]
The discriminants \(-3\) and \(-1/3\) are nonsquare units modulo
\(5\). Hence \(\Q_5[R]\) is the unramified quadratic extension
of \(\Q_5\), with ring of integers \(\Z_5[R]\).

An \(R\)-stable \(\Z_5\)-lattice in \(\Q_5^2\) is a
rank-one \(\Z_5[R]\)-module, hence a fractional ideal of this ring.
Since the extension is unramified, all these lattices are
homothetic under \(\Q_5^\times\). Moreover, \(\det R\) is a unit,
so a homothety class fixed by \(r\) has an \(R\)-stable
representative. Thus \(r\) fixes a unique homothety class of
lattices.

Conjugation by \(\Gamma\) sends \(r\) to \(r\) or \(r^{-1}\).
These elements have the same invariant lattice classes. It
follows that \(\Gamma\) preserves the unique class above.
For a representative \(\Lambda\),
\[
 \mathcal R=\End_{\Z_5}(\Lambda)
\]
is a \(\Gamma\)-stable maximal order. Every projective element
of \(\Gamma\) has a representative in
\(\GL(\Lambda)\), so it acts projectively on
\(\Lambda/5\Lambda\). This residual projective action is
irreducible, since the reduction of \(R\) has irreducible
characteristic polynomial.

Since \(A(\Q)[5]\neq0\), Lemma~\ref{lem:cyclotomic} gives
\(\Q(\zeta_5)\subseteq L\). Hence \(\Gamma\) has a cyclic quotient
of order \(4\). The orders of \(C_3\) and \(C_6\) are not divisible
by \(4\), while the abelianizations of \(D_6\) and \(D_{12}\) are
\(C_2\) and \(C_2\times C_2\), respectively. None has a cyclic
quotient of order \(4\), a contradiction.

The remaining groups with a rotation of order greater than \(2\)
are \(C_4\) and \(D_8\). For a rotation of order \(4\), the
trace-\(1\) representative has characteristic polynomial
\[
 X^2-X+\frac12.
\]
Its discriminant is \(-1\), a nonsquare unit modulo \(7\).
The same lattice argument, now over \(\Q_7\), produces a stable
maximal order with irreducible residual projective action.
Since \(A(\Q)[7]\neq0\), Lemma~\ref{lem:cyclotomic} gives
\(\Q(\zeta_7)\subseteq L\). Its degree is \(6\), which divides
neither \(4\) nor \(8\), again a contradiction.

We are left with
\[
 \Gamma\in\{C_1,C_2,D_4\}.
\]
In particular, \([L:\Q]\in\{1,2,4\}\).

\section{Completion of the proof}

Every geometric elliptic subvariety of \(A\) is defined over \(L\),
since it is the image of an integral multiple of an idempotent in
\(D\). Choose one, \(E\), and put \(F=A_L/E\). The elliptic curves
\(E\) and \(F\) are \(L\)-isogenous: the geometric rational
homomorphisms between the factors come from the matrix algebra
\(D\), all of whose endomorphisms are defined over \(L\).

For each \(p\in\{5,7\}\), the exact sequence
\[
 0\longrightarrow E\longrightarrow A_L\longrightarrow F
 \longrightarrow0
\]
shows that either \(E(L)\) or \(F(L)\) contains a point of order
\(p\). Thus their common \(L\)-isogeny class contains curves with
rational points of orders \(5\) and \(7\).
By Lemma~\ref{lem:combine}, it contains an elliptic curve over
\(L\) with a point of order \(35\).
Since \([L:\Q]\in\{1,2,4\}\), this contradicts the quartic
torsion classification. This proves Theorem~\ref{thm:main}.
\hfill\(\square\)

\section{Consequences for the census}

The groups
\href{https://genus-2-torsion.github.io/group.html?g=35}{\([35]\)}
and
\href{https://genus-2-torsion.github.io/group.html?g=70}{\([70]\)}
are impossible in the class ``geometrically split, simple over
\(\Q\)''. The theorem places no restriction on the geometric
splitting degree or on the integral geometric endomorphism ring.
It applies to abelian surfaces without assuming a principal
polarization.

\begin{thebibliography}{9}
\bibitem{DN}
M. Derickx and F. Najman,
\emph{Classification of torsion of elliptic curves over quartic
fields}, Theorem 1.1.
\href{https://arxiv.org/html/2412.16016v2}{arXiv:2412.16016v2} (2025).

\bibitem{Clark}
P. L. Clark,
\emph{Bounds for torsion on abelian varieties with integral moduli},
Theorem 5.
\href{https://arxiv.org/pdf/math/0407264}{arXiv:math/0407264v2} (2004).
\end{thebibliography}

\end{document}
